% GENERATED FILE. Edit canonical lessons, then run npm run generate:books.
% canonical-source-hash: fnv1a64:1997a8f50d928737
% canonical-lessons: PT-C475-nervosismo, PT-C475-irritar, PT-C475-impaciente, PT-C475-acalmar, PT-C475-alivio

\chapter{Nervos e calma}
\label{ch:pt-nervos-e-calma}
\hlchaptermodality{475}

\begin{chapteropening}
\textbf{By the end of this chapter:} \emph{I can talk about nerves, irritation and calming down.}

Complete the last lesson of chapter 475: I can talk about nerves, irritation and calming down.
\end{chapteropening}

\section[o nervosismo]{o nervosismo — nervousness}

\label{lesson:PT-C475-nervosismo}



\begin{quote}
Before the new one: say the Portuguese for tantrum, then the Portuguese for to whine.
\end{quote}

\subsection*{You'll want to know: o nervosismo}
\textbf{o nervosismo} — \textquotedblleft{}nervousness\textquotedblright{}.

\begin{quote}
Antes do exame, o nervosismo não me deixou dormir.
\end{quote}

\begin{quote}
\emph{Before the exam, nerves kept me awake.}
\end{quote}

Built on \emph{nervoso}.

\begin{cousinweb}
\emph{nervosismo} is built on \emph{nervoso} \textquotedblleft{}nervous\textquotedblright{}, from Latin \emph{nervōsus} \textquotedblleft{}sinewy, vigorous\textquotedblright{} (Lewis \& Short), from \emph{nervus} \textquotedblleft{}a sinew\textquotedblright{}.
\end{cousinweb}

\begin{grammarlens}[title={What you've built}]
One more word for telling, explaining and describing at B1.
\end{grammarlens}

\subsection*{Your turn}
\begin{itemize}
\raggedright
  \item \emph{Say it:} \emph{o nervosismo}
  \item \emph{Say it:} \emph{o nervosismo}, once more
  \item \emph{Say it:} a sentence of your own with \emph{o nervosismo}
  \item \emph{Recall:} say the Portuguese for tantrum, then the Portuguese for to whine, then say \emph{o nervosismo} again
\end{itemize}

\begin{tcolorbox}[breakable,colback=teal!4,colframe=teal!35!black,title={Before you move on}]
What does o nervosismo mean? (\textquotedblleft{}Nervousness\textquotedblright{}.) Say it once more.
\end{tcolorbox}
\section[irritar]{irritar — to irritate, to annoy}

\label{lesson:PT-C475-irritar}



\begin{quote}
Before the new one: say the Portuguese for to whine, then the Portuguese for nervousness.
\end{quote}

\subsection*{You'll want to know: irritar}
\textbf{irritar} — \textquotedblleft{}to irritate, to annoy\textquotedblright{}.

\begin{quote}
O barulho da rua irrita-me.
\end{quote}

\begin{quote}
\emph{The noise from the street annoys me.}
\end{quote}

\emph{Irritar-se} — to get annoyed: \emph{irrito-me com facilidade}.

\begin{cousinweb}
\emph{irritar} is a learned borrowing from Latin \emph{irrītāre} \textquotedblleft{}to provoke, to stir up\textquotedblright{} (EtymDB).
\end{cousinweb}

\begin{grammarlens}[title={What you've built}]
One more word for telling, explaining and describing at B1.
\end{grammarlens}

\subsection*{Your turn}
\begin{itemize}
\raggedright
  \item \emph{Say it:} \emph{irritar}
  \item \emph{Say it:} \emph{irritar}, once more
  \item \emph{Say it:} a sentence of your own with \emph{irritar}
  \item \emph{Recall:} say the Portuguese for to whine, then the Portuguese for nervousness, then say \emph{irritar} again
\end{itemize}

\begin{tcolorbox}[breakable,colback=teal!4,colframe=teal!35!black,title={Before you move on}]
What does irritar mean? (\textquotedblleft{}To irritate, to annoy\textquotedblright{}.) Say it once more.
\end{tcolorbox}
\section[impaciente]{impaciente — impatient}

\label{lesson:PT-C475-impaciente}



\begin{quote}
Before the new one: say the Portuguese for nervousness, then the Portuguese for to irritate.
\end{quote}

\subsection*{You'll want to know: impaciente}
\textbf{impaciente} — \textquotedblleft{}impatient\textquotedblright{}.

\begin{quote}
Fico impaciente quando o autocarro se atrasa.
\end{quote}

\begin{quote}
\emph{I get impatient when the bus is late.}
\end{quote}

The same form for a man and a woman.

\begin{cousinweb}
\emph{impaciente} is a learned borrowing from Latin \emph{impatiēns} \textquotedblleft{}unable to bear\textquotedblright{} (EtymDB), built on \emph{patī} \textquotedblleft{}to suffer, to bear\textquotedblright{}.
\end{cousinweb}

\begin{grammarlens}[title={What you've built}]
One more word for telling, explaining and describing at B1.
\end{grammarlens}

\subsection*{Your turn}
\begin{itemize}
\raggedright
  \item \emph{Say it:} \emph{impaciente}
  \item \emph{Say it:} \emph{impaciente}, once more
  \item \emph{Say it:} a sentence of your own with \emph{impaciente}
  \item \emph{Recall:} say the Portuguese for nervousness, then the Portuguese for to irritate, then say \emph{impaciente} again
\end{itemize}

\begin{tcolorbox}[breakable,colback=teal!4,colframe=teal!35!black,title={Before you move on}]
What does impaciente mean? (\textquotedblleft{}Impatient\textquotedblright{}.) Say it once more.
\end{tcolorbox}
\section[acalmar]{acalmar — to calm (down)}

\label{lesson:PT-C475-acalmar}



\begin{quote}
Before the new one: say the Portuguese for to irritate, then the Portuguese for impatient.
\end{quote}

\subsection*{You'll want to know: acalmar}
\textbf{acalmar} — \textquotedblleft{}to calm (down)\textquotedblright{}.

\begin{quote}
A mãe acalmou o bebé com uma canção.
\end{quote}

\begin{quote}
\emph{The mother calmed the baby with a song.}
\end{quote}

\emph{Acalmar-se} — to calm yourself down: \emph{acalmei-me}.

\begin{cousinweb}
\emph{acalmar} is built on \emph{calma} \textquotedblleft{}calm\textquotedblright{} (EtymDB), from Late Latin \emph{cauma} \textquotedblleft{}the heat of the midday sun\textquotedblright{}, from Greek \emph{kaûma}: at noon everything goes still.
\end{cousinweb}

\begin{grammarlens}[title={What you've built}]
One more word for telling, explaining and describing at B1.
\end{grammarlens}

\subsection*{Your turn}
\begin{itemize}
\raggedright
  \item \emph{Say it:} \emph{acalmar}
  \item \emph{Say it:} \emph{acalmar}, once more
  \item \emph{Say it:} a sentence of your own with \emph{acalmar}
  \item \emph{Recall:} say the Portuguese for to irritate, then the Portuguese for impatient, then say \emph{acalmar} again
\end{itemize}

\begin{tcolorbox}[breakable,colback=teal!4,colframe=teal!35!black,title={Before you move on}]
What does acalmar mean? (\textquotedblleft{}To calm (down)\textquotedblright{}.) Say it once more.
\end{tcolorbox}
\section[o alívio]{o alívio — relief}

\label{lesson:PT-C475-alivio}



\begin{quote}
Before the new one: say the Portuguese for impatient, then the Portuguese for to calm.
\end{quote}

\subsection*{You'll want to know: o alívio}
\textbf{o alívio} — \textquotedblleft{}relief\textquotedblright{}.

\begin{quote}
Que alívio! O exame correu bem.
\end{quote}

\begin{quote}
\emph{What a relief! The exam went well.}
\end{quote}

\emph{Sentir alívio} — to feel relieved.

\begin{cousinweb}
\emph{alívio} comes from \emph{aliviar} \textquotedblleft{}to relieve\textquotedblright{} (EtymDB), which is inherited from Latin \emph{alleviāre} \textquotedblleft{}to lighten\textquotedblright{}, built on \emph{levis} \textquotedblleft{}light\textquotedblright{}.
\end{cousinweb}

\begin{grammarlens}[title={What you've built}]
One more word for telling, explaining and describing at B1.
\end{grammarlens}

\subsection*{Your turn}
\begin{itemize}
\raggedright
  \item \emph{Say it:} \emph{o alívio}
  \item \emph{Say it:} \emph{o alívio}, once more
  \item \emph{Say it:} a sentence of your own with \emph{o alívio}
  \item \emph{Recall:} say the Portuguese for impatient, then the Portuguese for to calm, then say \emph{o alívio} again
\end{itemize}

\begin{tcolorbox}[breakable,colback=teal!4,colframe=teal!35!black,title={Before you move on}]
What does o alívio mean? (\textquotedblleft{}Relief\textquotedblright{}.) Say it once more.
\end{tcolorbox}
